\section{Reconhecimento Óptico de Caracteres (OCR)}
\label{sec:ocr}


\paragraph{Modelagem de contexto.}
Tradicionalmente usam-se LSTMs/GRUs bidirecionais; 
em \textit{Edge AI}, porém, LSTMs impõem gargalos de memória e restrições 
em \textit{runtimes} C++, sendo por vezes substituídas por convoluções 
temporais 1D de núcleos largos que emulam o campo receptivo contextual 
sem o custo da recorrência \cite{shi2017crnn}.  


\subsection{Edge OCR, Quantização e TinyML}
A transição de modelos de OCR de servidores em nuvem para microcontroladores com
SRAM restrita impõe o paradigma do TinyML e da Computação de Borda \cite{ademola2023resource}.

\paragraph{Quantization-Aware Training (QAT).}
